\section{Introduction}
\label{sec:intro}

A \emph{succinct argument for machine execution} convinces a verifier that a program, run on a given instruction-set architecture, produced a given output, with a proof far shorter and a check far cheaper than the execution. Such systems are commonly called zero-knowledge virtual machines (\zkVM{}s), although the proofs studied here are succinct but not hiding. The benchmark throughout is an Ethereum block of several hundred million instructions~\citep{ethproofs}, a hash-intensive workload that we use as a benchmark rather than as the application.

\subsection{Succinct arguments for machine execution}
\label{sec:intro-background}

A \zkVM{} fixes an instruction set, a register file and a memory model, and proves that a program for that machine ran correctly. Representative systems that prove RISC-V programs are SP1, OpenVM, ZisK and Jolt~\citep{sp1,openvm,zisk,jolt}, and system-level comparisons are restricted to these four. Such a system has two halves, a front end and a back end. The \emph{arithmetisation} turns an execution into tables whose rows satisfy polynomial constraints and between which multiset relations hold; the \emph{commitment} and its proximity test let the prover show that it knows such tables. The first decides how many field elements an executed instruction becomes, and the second is where the prover pays for each.

\paragraph{Motivation.} \Ziren{} is a production \zkVM{} for MIPS32. It proves 77 user-mode integer instructions of MIPS32r2 with the branch-delay slot (Appendix~\ref{sec:app-isa}, which also names the few it leaves out), and it has proved Ethereum mainnet blocks end to end in production; earlier proof systems address MIPS-\emph{like} models or offer MIPS as one mode among several (\S\ref{sec:related}). The specification is \emph{stable}: its owner moved its roadmap to RISC-V in 2021~\citep{mips-riscv-2021,mips-at-40}, so the instruction set, and the verifying keys derived from it, do not change. It is \emph{uniform}: fixed 32-bit encodings in a handful of formats let one row carry an instruction's whole frame. And its toolchains are \emph{mature}: Rust, Go and C programs compile to the same little-endian ELF, and the architecture is widely deployed~\citep{mips-shipped}. The fault-proof software of Ethereum rollups, Optimism's Cannon and Kona~\citep{cannon,kona}, is therefore in range, subject to compatible runtime conventions and image formats.

\paragraph{Challenges.} Four properties shape the design. The instruction after a branch executes before control transfers,\footnote{For jumps and conditional branches alike, and after a conditional branch whether or not it is taken~\citep{mips32-isa}; an argument system must constrain the delay slot, since normalising it away in the emulator proves a different machine.} so an instruction's successor is data rather than the following address. The relations that hold the machine together are multiset relations \emph{across} tables, so discharging them table by table would cost one argument per bus per table. A shard of the block workload commits up to about $3.6\cdot10^4$ columns whose heights run from 32 rows to millions, so neither one commitment per column nor a common height is affordable. And a 32-bit word does not fit the 31-bit field, so every word is a vector of bytes.

\subsection{Overview of Ziren}
\label{sec:intro-approach}

In \Ziren{}, every executed instruction is one row of the table (the \emph{chip}) for its opcode, and there is no central CPU table. Each chip embeds what we call an \emph{instruction frame}, the fixed group of columns and bus interactions that surrounds every row's own logic and makes the row an executed instruction whatever it computes: the frame fetches the row's instruction from the committed program, performs its register and memory accesses at fixed sub-cycle positions, and receives and sends the machine state on a bus; the state is a program-counter pair, which absorbs the delay slot (\S\ref{sec:air-frame}). The committed trace is thus the execution regrouped by opcode, the device of vRAM~\citep{vram}, with the state bus as its permutation check.

Above the rows there is one argument per shard. Registers and memory share one offline memory-checking argument, carried across shards by an elliptic-curve digest; every bus is discharged by one \LogUp{} instance and every constraint by one zerocheck; and the jagged sumcheck of Hemo et al.~\citep{jagged-pcs} reduces the claims on tens of thousands of columns of differing heights to one evaluation, which a single batched \WHIR{} opening~\citep{whir} authenticates against every commitment. Recursion composes the shard proofs under a verifying-key allowlist. A block proof is
\[
  \mathsf{Rec}\Big[\,\mathsf{FS}\big[\,\mathsf{JaggedWHIR}\big[\,\mathsf{LogUpGKR} \wedge \mathsf{Zerocheck}\,\big[\mathcal R_{\mathrm{shard}}\big]\big]\big]\Big],
\]
and Figure~\ref{fig:architecture} draws the pipeline. Section~\ref{sec:architecture} describes the system as a stack of these arguments ending in a compressed proof for a native verifier, and two applications built on it, which verify that proof inside other proof systems for on-chain and garbled verifiers (\S\ref{sec:applications}).

\paragraph{Contributions.} The commitment scheme, proximity test, lookup argument and recursion shape \Ziren{} builds on are prior work, used as published; our contributions are the machine and three techniques around it.
\begin{itemize}
\item \emph{The machine.} An arithmetisation of MIPS32 with no CPU table, in which an instruction frame carries the delay slot as a program-counter pair (Sections~\ref{sec:isa} and~\ref{sec:air}); the concrete composition of the jagged reduction with \WHIR{}, with its multi-root layout, parameter schedule and recursive verifier (Sections~\ref{sec:pcs} and~\ref{sec:iopp}); and a verifying-key allowlist of about $6.4\cdot10^{3}$ keys enumerated from the machine's declared shapes without executing a guest (Section~\ref{sec:recursion}).
\item \emph{Determinism in \Lean{}.} The constraint system of every chip the extraction models (57 chips) is taken from the Rust description the prover evaluates, and its determinism is stated as a \Lean{} theorem; a propagation analysis derives each output from the inputs, and the generator replays that derivation as a proof, one step per determined column, over a library of hand-proved gadget lemmas (\S\ref{sec:verification}). All \ndettotal{} theorems are closed, with no axiom beyond \Lean{}'s standard three.
\item \emph{Pipelined proving.} Recursion overlaps core proving: leaves and compose nodes are proved on the GPUs while the block's core shards are still being proved, so only two serial composes remain after the last shard proof, and trace generation and the proof-of-work searches run on the device (Section~\ref{sec:perf-pipeline}); its effect on wall-clock time is not separated from the rest (\S\ref{sec:perf-threats}).
\item \emph{Four options for cross-shard memory.} Multiset hashes over an elliptic curve and over a lattice, a global \LogUp{} challenge drawn from a hash chain over every shard's commitment, and Merkle roots over touched memory (\S\ref{sec:air-global}). The elliptic-curve hash is the one component that is not hash-based; against it, on the same GPUs, the global challenge is $29$ to $42\,\%$ slower and Merkle roots are $16.6\,\%$ faster on one GPU and $3.4\,\%$ slower on four, and the lattice hash, costed, would need at least $19.6\times$ today's committed area. Section~\ref{sec:pq} gives what each rests on against a quantum adversary.
\end{itemize}

\subsection{Costs}
\label{sec:intro-costs}

\paragraph{Cost model.} \Ziren{} commits Merkle trees over Reed--Solomon encodings of a 31-bit field, where every cell costs the same whatever it holds; a small value is no cheaper than a large one, unlike under a commitment built from multi-scalar multiplication, where small values are cheaper~\citep{jolt}. The unit of cost is therefore the \emph{cell}, together with the bus interactions each row raises, which size the lookup argument. Cost has two independent parts. \emph{Instruction efficiency} is the cells and interactions per executed instruction, a property of the arithmetisation that does not change with the hardware; \emph{proving efficiency} is the rate at which an implementation turns cells into a proof, in guest cycles per second on named hardware. End-to-end time is the work the first gives an execution, divided by the rate the second gives, and Section~\ref{sec:performance} states which of the two each measurement reports.

\paragraph{Prover costs.} An executed instruction pays for its frame before it computes anything (\S\ref{sec:air-cost}), and the operation adds from four columns (\texttt{AddSub}) to over a hundred (\texttt{DivRem}) among the arithmetic chips (\S\S\ref{sec:base} and~\ref{sec:ext}); block 25\,955\,640 commits about 59 cells and raises about 26 bus interactions per executed instruction. Cells per instruction is a property of the workload, not a constant: a precompile call is one guest cycle but many rows, and a shard closes on its committed area, so guest cycles, the unit block-proving dashboards report, do not measure prover work (\S\ref{sec:perf-estimate}). On 16 consecutive mainnet blocks proved with both Reth and Geth, one NVIDIA RTX~5090 proves at up to $10.3$\,MHz and eight at up to $80$\,MHz, yet the Reth guest, at lower throughput, proves each block $3.1$--$7.0$ times faster than Geth on one GPU (Table~\ref{tab:clients}).

\paragraph{Verifier costs.} The verifier reads one compressed proof of 274\,KiB and checks it in about 73\,ms on one host; the size does not depend on the execution's length and is set by the opening's schedule (\S\ref{sec:perf-proofsize}). Every proof stage is above 100 bits, and the \emph{composite security} of a block, the level of its whole recursion tree with every node's error summed, is $93.5$ bits for the measured block's tree of about 150 proofs. These figures bound the \emph{interactive} protocol; \S\ref{sec:iopp-soundness} gives the union-bound accounting and the level after Fiat--Shamir compilation.

\subsection{Auditable soundness}
\label{sec:intro-audit}

An accepting proof implies a satisfying trace; that the trace is an execution of the advertised machine needs the arithmetisation to admit one behaviour per input, which the proof system cannot establish about itself. \Ziren{} audits its per-opcode constraints mechanically: it extracts the determinism of every chip it models as a theorem statement from the constraint description and discharges all \ndettotal{} of them in \Lean{}. A lookup-only design instead relies on tables fixed by the instruction set, which are easier to audit~\citep{jolt}.

\begin{theorem}[Conditional argument statement, informal]
\label{thm:main-informal}
Let $\mathsf{Emulate}$ be the MIPS32 execution relation of Section~\ref{sec:isa}, and let the block proof be the protocol of Section~\ref{sec:architecture} over the 31-bit field \KB{}. Honest executions that halt with exit code zero produce accepting proofs (\S\ref{sec:statement}). Suppose that every node's protocol is round-by-round knowledge sound, that child extractors accept the transcripts produced by parent extractors, that the recursion programs enforce the compose relation, and that Assumption~\ref{ass:digest} holds (\S\ref{sec:security-relation}). Then the recursion tree is knowledge sound for its satisfying witness rows in the random-oracle model, with the conditional error bound of Theorem~\ref{thm:composition}; \S\ref{sec:iopp-layers} separates this statement into a functional IOP, an interactive commit-and-open argument and its Fiat--Shamir compilation. Identifying those rows with a unique $\mathsf{Emulate}$ execution additionally requires the premises of Theorem~\ref{thm:unique}, chip determinism, proved in \Lean{} for every extracted chip (\S\ref{sec:verification}), and real-row exhaustiveness (\S\ref{sec:security-relation}).
\end{theorem}

\begin{theorem}[Chip determinism, informal]
\label{thm:det-informal}
For every chip $U$ of the core machine that the extraction models, there is a mechanically extracted \Lean{} theorem stating that two rows of $U$ that satisfy its constraints and agree on their inputs (their bus inputs and the values the constraints leave to the prover by design) agree on their bus outputs. All \ndettotal{} such theorems are closed (\S\ref{sec:verification}). Together with the bus arguments, chip determinism implies that an accepting shard proof determines, up to permutation, every row reachable from the entry state along the state and memory buses, given the initial values of the words the shard touches and those declared free values (Theorem~\ref{thm:unique}); that this exhausts the real rows is argued, not proved (\S\ref{sec:scope}).
\end{theorem}

\subsection{Technical overview}
\label{sec:techniques}

We illustrate the arithmetisation on a single instruction. Consider a MIPS32 \texttt{and} executed at clock $\mathsf{clk}$. It occupies one row of the \texttt{Bitwise} chip, which is 37 columns wide. Thirty columns form the instruction frame (\S\ref{sec:air-cost}), and two further columns hold the program-counter pair. Through the frame, the row sends $(\mathsf{pc}, \mathsf{opcode}, \ldots)$ on the program bus; on the state bus it receives $(\mathsf{shard}, \mathsf{clk}, \mathsf{pc}, \mathsf{next\_pc})$ and sends $(\mathsf{shard}, \mathsf{clk}+5, \mathsf{next\_pc}, \mathsf{next\_pc}+4)$; and it performs its three register accesses on the memory bus. The remaining five columns are four selectors and a gate. No polynomial constraint computes the result: it is established by four lookups into a $2^{16}$-entry byte table, the decomposition of the original Jolt design (\S\ref{sec:air-bytes}).

Besides the delay slot (\S\ref{sec:air-frame}), the general case raises two further difficulties: a limb identity that holds in $\F$ need not hold over the integers; and memory accesses cross shard boundaries (\S\ref{sec:air-global}). Section~\ref{sec:argument} states the protocol that operates on these rows.

\subsection{Related work}
\label{sec:related}

Table~\ref{tab:related} places the systems discussed along the axes of this introduction. Its axes are design axes: no cross-system performance claim follows from them, and no same-hardware comparison is attempted. Machine-checked constraint-system correctness is compared only with the published verification work cited below.

\begin{table}[t]
\centering
\scriptsize\setlength{\tabcolsep}{3pt}
\caption{Design axes of \Ziren{} and four systems it is most often read against, from their public descriptions at the release named in the last column, retrieved on 2026-09-24. The \emph{Regime} column names the basis on which each project states its soundness level; regimes differ and levels are not comparable without the accounting of Section~\ref{sec:iopp-soundness}.}
\label{tab:related}
\begin{tabular}{@{}>{\raggedright\arraybackslash}p{0.125\textwidth}
                  >{\raggedright\arraybackslash}p{0.095\textwidth}
                  >{\raggedright\arraybackslash}p{0.090\textwidth}
                  >{\raggedright\arraybackslash}p{0.105\textwidth}
                  >{\raggedright\arraybackslash}p{0.085\textwidth}
                  >{\raggedright\arraybackslash}p{0.120\textwidth}
                  >{\raggedright\arraybackslash}p{0.135\textwidth}
                  >{\raggedright\arraybackslash}p{0.075\textwidth}@{}}
\toprule
System & ISA & Field & Instruction tables & Lookup & Commitment / IOPP & Regime & Release \\
\midrule
\Ziren{} & MIPS32 & KoalaBear ($2^{31}$) & chip per opcode, frames & \LogUp{} & Jagged over \WHIR{} & unique decoding, Johnson bound at the root; union bound & this paper \\
SP1 Hypercube~\citep{sp1-hypercube} & RV64IM & KoalaBear & chip per opcode, readers & \LogUp{} & Jagged over BaseFold & unique decoding, no proximity-gap conjecture & v6.8.0 \\
ZisK~\citep{zisk} & RV64IMA & Goldilocks & micro-operations, PIL2 & LogUp & PIL2 (e)STARK & not stated & v1.3.0-alpha \\
OpenVM\newline\citep{openvm} & own ISA, RV32IM transpiler & BabyBear & executor chips, adapter and core & LogUp & SWIRL: stacked \WHIR{} & unique decoding low in the tree, list decoding above & v2.0.2 \\
Jolt~\citep{jolt} & RV64IMAC & BN254 scalar (Dory), 128-bit (Akita) & Twist and Shout, small R1CS & Twist and Shout & Dory (default), Akita (lattice) & not stated & v0.3.0-alpha \\
\bottomrule
\end{tabular}
\end{table}

\paragraph{Chip-decomposed machines.} SP1~\citep{sp1} decomposes a RISC-V machine into one chip per opcode with \LogUp{} and recursive aggregation, SP1 Hypercube~\citep{sp1-hypercube} adds the jagged commitment with a shard-level lookup and zerocheck, and SP1 also has a just-in-time executor~\citep{sp1-jit}. SP1 has no CPU table either. Two of SP1's devices are used here and credited where they appear: the elliptic-curve cross-shard digest~\citep{sp1-turbo-memo} and the register shadow read. \Ziren{} differs in the instruction set, in an instruction frame that carries the delay slot and replaces the CPU table, in its arithmetisation of MIPS32 family by family, in opening the jagged commitment with \WHIR{} rather than BaseFold, and in a union-bound soundness level.

SP1 counts query soundness in the unique-decoding regime and relies on proximity-gap theorems rather than conjectures~\citep{sp1-security}; \Ziren{} does the same below the root and uses the Johnson bound, also proved, at the root (\S\ref{sec:iopp-soundness}). OpenVM~\citep{openvm} likewise has no CPU chip; the instruction frame plays the role of SP1's state-and-reader columns and of OpenVM's adapter for MIPS32, and Section~\ref{sec:air-frame} compares the three.

\paragraph{The commitment layer.} OpenVM's SWIRL combines stacked \WHIR{} with \LogUp{}~\citep{swirl}, the same pair of ideas as our commitment layer. It pads columns to power-of-two heights with a univariate skip and reduces openings by a batch sumcheck fused with the lookup's input layer, where we round heights to a multiple of 32 and use the jagged sumcheck. Three further differences follow: our \WHIR{} rounds fold several variables at a time; \KB{} rather than BabyBear changes the challenge extension and hence the fingerprinting term; and SWIRL composes round-by-round errors by their maximum per layer, a basis not comparable with our union bound. Ceno's implementation offers a jagged reduction over BaseFold by default, and BaseFold or \WHIR{} alone as alternatives~\citep{ceno-repo}, while the Ceno paper leaves the commitment abstract~\citep{ceno}; the jagged paper itself anticipates \WHIR{} in place of its FRI-based opening~\citep[Rem.~7.1]{jagged-pcs}. We claim the concrete composition and its soundness accounting, not the idea of combining them.

\paragraph{Where the prover's work is put.} Jolt~\citep{jolt} evaluates instructions by lookups rather than per-opcode constraints; since v0.2.0 it uses Twist and Shout~\citep{twist-shout} in place of the Lasso construction~\citep{lasso} of its original design, with a small R1CS instance for instruction fetch. \Ziren{} uses lookups for byte-granular facts and buses and evaluates instructions by constraints. The two meet at three points: the byte table, the original Jolt design's decomposition of an instruction table (\S\ref{sec:air-bytes}), since made unnecessary by Shout~\citep{twist-shout}; the chip-per-opcode layout, the regrouping by instruction that Jolt, after vRAM~\citep{vram}, names as one way to route a step to its table (\S\ref{sec:air-frame}); and the timestamped memory check the original Jolt design took from Spice~\citep{spice}, where Jolt now uses Twist (\S\ref{sec:air-memory}). ZisK~\citep{zisk} instead lowers RV64IMA execution to micro-operations constrained in PIL2 and commits a STARK trace. Section~\ref{sec:performance} measures this system against itself.

\paragraph{Verifiable RAM and regrouped traces.} Three of this system's devices come from one line of work: Blum et al.'s offline memory check~\citep{blum1994} with Spice's timestamps~\citep{spice}, which verifies a read--write memory without reordering the trace and which \Ziren{}'s registers and memory use (\S\ref{sec:prelim-lookups}); vRAM's regrouping of the trace by instruction, proved a permutation~\citep{vram}, which Ceno's per-opcode design~\citep{ceno} published in multiset form and which the frame realises with the delay slot's counter pair, a shard field and a proved chain lemma (\S\ref{sec:air-frame}); and Ceno's initialisation of only the memory a run touches, which \Ziren{} shares (\S\ref{sec:air-memory}). Ceno's basic-block design has no counterpart here, and its implementation is per opcode~\citep{ceno-repo}.

\paragraph{Other systems.} RISC Zero~\citep{risc0} proves one trace of machine cycles under a fixed constraint set, selected per cycle by public control columns, with a sorted-column permutation for memory and a DEEP-ALI and FRI STARK, and, unlike \Ziren{}, claims zero knowledge (2023 draft). Cairo~\citep{cairo} and Valida~\citep{valida} design an instruction set for proving instead of inheriting a hardware one. Nexus~1.0 folded executions with Nova-family schemes in a tree of proof-carrying data~\citep{nexus}; Nexus~3.0 replaced folding with a circle STARK over Mersenne-31, a single trace with a CPU component and a logarithmic-derivative memory check~\citep{nexus3}. Airbender~\citep{airbender} targets Mersenne-31 with a degree-two arithmetisation and DEEP-FRI, and Cysic~\citep{cysic} pursues dedicated proving hardware, outside the CPU--GPU scope of this paper. Among MIPS systems, Zilch's zMIPS verifies a MIPS-like instruction model and composes instruction proofs~\citep{zilch}; o1VM documents a segmented prover with MIPS and RISC-V modes~\citep{o1vm}; Cannon~\citep{cannon} is a single-step MIPS interpreter for interactive fault proofs whose suite is one of our conformance inputs, and Kona~\citep{kona} the fault-proof program of the same ecosystem. On constrained devices, proofs today serve authentication and attestation~\citep{zk-iot-survey,zpie,zekra,zra}; a \zkVM{} for the device's instruction set makes its software execution the statement instead.

\paragraph{Primitives and verification.} The field, $\Poseidon{}$~\citep{poseidon2}, constraint interfaces and transforms come from Plonky3~\citep{plonky3}. Circle STARKs~\citep{circle-starks} reach smooth Mersenne-31 domains through the circle group, and Binius~\citep{binius,fri-binius} avoids splitting words with binary towers; \KB{} uses ordinary Reed--Solomon domains. The proximity line runs from FRI~\citep{ben2018fast} through DEEP-FRI~\citep{deep-fri} and proximity gaps~\citep{proximity-gaps} to BaseFold~\citep{basefold}, STIR~\citep{stir} and \WHIR{}~\citep{whir}; Ligero~\citep{ligero} and Brakedown~\citep{brakedown} are linear-time alternatives. The jagged commitment is that of Hemo et al.~\citep{jagged-pcs}, with the branching-program evaluation of Holmgren and Rothblum~\citep{hr18}. The lookup is Hab\"ock's logarithmic derivative~\citep{haboeck} in its GKR realisation~\citep{logup-gkr}; Plookup~\citep{plookup} and cq~\citep{cq} are univariate alternatives that do not transfer to this multilinear many-bus workload. Picus~\citep{picus} introduced the determinism formulation and an SMT solver for it; we discharge it in \Lean{}, which gives kernel-checked proofs. OpenVM's verification project proves refinement of its RV32IM and precompile chips to Lean specifications~\citep{openvm-fv}, a stronger property than determinism, for a different ISA. The soundness accounting follows the \WHIR{} analysis~\citep{whir}, computed with \texttt{soundcalc}~\citep{soundcalc}, in the style of the ethSTARK documentation~\citep{ethstark}.

\subsection{Scope and limitations}
\label{sec:scope}

This paper measures and accounts only \Ziren{}, the path to the compressed proof; the applications built on that proof (\S\ref{sec:applications}) are outside its scope. Soundness rests on a computational assumption for the cross-shard digest in addition to the information-theoretic bounds of the proximity test, and that assumption, unlike the rest of the argument, does not survive a quantum adversary (\S\ref{sec:pq}). The determinism pipeline takes any chip written against the prover's constraint interface; we have run it on the core machine, and the two SHA-256 control chips and the recursion machine's chips, written against the same interface, are its next inputs rather than a new method. Real-row exhaustiveness and Theorem~\ref{thm:unique} itself are argued on paper (\S\ref{sec:security-relation}). The evaluation uses one workload class on one GPU model. The verifying-key allowlist is enumerated from the machine's declared chip-set clusters, which are curated from its target workloads rather than exhaustive, so a shard outside them is refused rather than admitted (\S\ref{sec:recursion}).
